\subsection{Negative quadratic drift away from the secondary fold}
\label{sec:direct-buffered-negative}

We first treat the part of
\eqref{eq:direct-negative-secondary-scales} on which
\begin{equation}
\label{eq:direct-buffered-region}
 b^2-4a\le-\nu<0 .
\end{equation}
The fixed margin implies \(a\ge\nu/4\). Both secondary passages therefore
have regular prefixes. The remaining comparison must match their exits
at the same central height, even when the individual actions become
large.

\begin{proposition}[Buffered negative quadratic drift]
\label{prop:direct-buffered-negative}
For every fixed \(\nu>0\), the return displacement \(\Delta_p\) on
\eqref{eq:direct-buffered-region} has at most
three distinct zeros, for all sufficiently small independent
\(s,r>0\) and all central heights in \(\cH_p\).
Moreover, the logarithmic slope ratio \(\mathcal L\) has
at most two distinct zeros on \(\cH_p\).
\end{proposition}

\subsubsection{Actions and the regular prefix}

On a small fixed interval \(0\le t\le v_1\),
\[
 D_\sigma(t)=ar^2-\sigma brt+t^2
                    -\sigma t^3R_q(\sigma st)
                 \asymp r^2+t^2.
\]
Indeed the two quadratic forms have a uniform positive
margin by \eqref{eq:direct-buffered-region}, and the cubic
term is absorbed by reducing \(v_1\).
Set
\begin{equation}
\label{eq:direct-buffered-actions}
 \mathcal A_\sigma(\rho)=
                    \int_0^\rho t/D_\sigma(t)\,dt .
\end{equation}
The positive action is defined up to any fixed positive primary radius;
we use the negative action first on compact subintervals below its
far root near \(\rho=1\).
At \(\rho_i=K_0r\), the central scale
\(\ell=sr,\ Q=s^3r^2\) has drifts
\[
 a-\sigma bx+x^2-\sigma rx^3R_q(\sigma srx),
                 \qquad 0\le x\le K_0 .
\]
They have a fixed positive margin for sufficiently small \(r\).
The central-height construction and composition with the
annular coordinate therefore give
\eqref{eq:direct-annulus-prefix}, with
\(\mathcal A_{\sigma,i}=\mathcal A_\sigma(K_0r)\).

At a common small primary radius,
\begin{equation}
\label{eq:direct-buffered-action-difference}
 |\mathcal A_+(\rho)-\mathcal A_-(\rho)|
 \le C\int_0^\rho
          \frac{brt^2+t^4}{(r^2+t^2)^2}\,dt
 \le C(b+\rho).
\end{equation}
The substitution \(t=rx\) gives an integrable density for the first
term, and the second integrand is bounded. This bounds the difference
even when the individual actions grow like \(\log(1/r)\).

\subsubsection{The intermediate annulus}

We prove that there are fixed \(M_0>2K_0\) and \(v_0>0\) such that the
interval of central heights in \(\mathcal H_p\) with positive thin exit in
\begin{equation}
\label{eq:direct-buffered-intermediate}
 M_0r\le\rho\le v_0
\end{equation}
contains at most three zeros of \(\Delta_p\).
We first construct the negative exit for each of these central heights.
It has the form \(\omega=\Omega\rho\), where
\begin{equation}
\label{eq:direct-buffered-comparable}
 0<\Omega_0\le\Omega\le\Omega_1,\qquad
 |\log\Omega|\le C(b+\rho+s).
\end{equation}
By \eqref{eq:direct-negative-secondary-scales}, \(0\le b\le1\).
Fix the lower action-slope
margin \(c_{\rm a}>0\) on the whole regular range
\(K_0r\le t\le v_1\).  Choose fixed factors
\(0<\Omega_0<1<\Omega_1\) so far apart that
\begin{equation}
\label{eq:direct-buffered-trial-separation}
 c_{\rm a}\min\{|\log\Omega_0|,\log\Omega_1\}>8C_{\rm act},
\end{equation}
where \(C_{\rm act}\) is a fixed bound for the \(b\)-dependent part of the
full \(\Phi(1;\rho)\) comparison, including the action and bounded
height-factor terms, when \(0\le b\le1\) and \(0<\rho\le1\).
The terms proportional to \(s\log\Omega\) are handled after the fixed
trial factors have been chosen.
Then choose \(M_0>2K_0/\Omega_0\) and \(v_0>0\) so small
that \(\Omega_1v_0<v_1\).  On
\([\Omega_0\rho,\Omega_1\rho]\), with \(M_0r\le\rho\le v_0\),
the positive margin in \eqref{eq:direct-buffered-region} gives
\begin{equation}
\label{eq:direct-buffered-negative-action-slope}
 0<c_{\rm a}\le
 \frac{t^2}{D_-(t)}\le C_{\rm a},\qquad
 \left|\frac{t^2}{D_-(t)}-\frac{t^2}{D_+(t)}\right|\le C(b+t),
\end{equation}
after reducing \(v_0\).  Let \(h_\sigma(\rho)\) denote the central height
obtained from \eqref{eq:direct-annulus-source} at the first thin exit
\(\rho\) on branch \(\sigma\).  The event equation and
\eqref{eq:direct-annulus-source-rate} imply that \(h_\sigma\) is \(C^3\)
and strictly decreasing, and that
\begin{equation}
\label{eq:direct-buffered-height-comparison}
 \Phi(\Omega;\rho):=s\log\frac{h_-(\Omega\rho)}{h_+(\rho)},\qquad
 \partial_{\log\Omega}\Phi
 =-\frac{(\Omega\rho)^2}{D_-(\Omega\rho)}+O(s)\le-\frac12c_{\rm a}.
\end{equation}
At \(\Omega=1\), the action difference estimate
\eqref{eq:direct-buffered-action-difference}, together with the bounded
height factors in \eqref{eq:direct-annulus-source}, gives
\begin{equation}
\label{eq:direct-buffered-height-at-one}
 |\Phi(1;\rho)|\le C_{\rm act}+C(v_0+s),
\end{equation}
with \(C_{\rm act}\) as chosen above. First fix the trial separation in
\eqref{eq:direct-buffered-trial-separation}; then reduce \(v_0\) and
finally \(s\).  The bound in
\eqref{eq:direct-buffered-height-at-one} is then smaller than
\(\frac14c_{\rm a}\min\{|\log\Omega_0|,\log\Omega_1\}\) uniformly in
\(b\in[0,1]\).  Integrating
\eqref{eq:direct-buffered-height-comparison} then gives
\(\Phi(\Omega_0;\rho)>0>\Phi(\Omega_1;\rho)\), uniformly for
\(M_0r\le\rho\le v_0\). The intermediate value theorem therefore gives an exit
\(\omega=\Omega\rho\) with \(h_-(\omega)=h_+(\rho)\).  Strict decrease in
\eqref{eq:direct-buffered-height-comparison} gives uniqueness of the
matching radius, and each passage has a unique first exit by
\eqref{eq:direct-annulus-source-rate}.

At the matching radius, combine the equal-height equation with
\eqref{eq:direct-buffered-action-difference}. The lower bound
in \eqref{eq:direct-buffered-height-comparison} yields
\begin{equation}
\label{eq:direct-buffered-log-omega-bound}
 |\log\Omega|\le C(b+\rho+s),
\end{equation}
after absorbing the \(O(s)\) height factors, proving the second bound in
\eqref{eq:direct-buffered-comparable} for sufficiently small parameters.
Differentiation of the equation of equal central heights gives
\[
 \mathcal E\log\omega=G_+(\rho)/G_-(\omega),\qquad
 \mathcal E\log\Omega=O(u+v).
\]

For \(b\ge b_{*,-}>0\), where the threshold will be fixed by the
small-\(b\) estimate below, the preceding prefixes and comparable exit
radii meet \cref{cor:direct-annular-pair} after choosing the annular
margins small relative to \(b_{*,-}\), then reducing \(s\). It gives
\(\mathcal L''<0\) in \(T=-s\log h\), and at most three zeros of \(\Delta_p\).

For small \(b\), curvature need not have a uniform margin. We bound the
endpoint slope quotient and use the exterior contraction. The quotient is
\[
 \frac{I_{-,h}}{I_{+,h}}
       =\frac{M_-(\omega)}{M_+(\rho)}
                         \frac{G_+(\rho)}{G_-(\omega)}.
\]
The logarithm of this quotient has, by
\eqref{eq:direct-annulus-slope-calibration},
\eqref{eq:direct-buffered-comparable}, and
\eqref{eq:direct-annulus-source-rate}, the upper bound
\[
 C\{b+u_0+v_0+s(1+|\log s|)\}
       +2C_{\rm cone}\sqrt{z_*}.
\]
Choose \(b_{*,-}>0\) and the annular margins small to make this less than
\(-\log\Lambda_+\) for \(b\le b_{*,-}\), after reducing \(s\).
Then \(\Pi_p'I_{-,h}/I_{+,h}<1\), so \(\Delta_p\) has at most
one zero. Shrink the annular margins further for the curvature
argument with this same \(b_{*,-}\), then reduce \(s\).
This proves the three-zero bound on
\eqref{eq:direct-buffered-intermediate}.

\subsubsection{Exits on a compact primary interval}
\label{sec:direct-buffered-compact-primary}

Fix \(0<\rho_{\rm lo}<\rho_{\rm hi}<\infty\). Suppose the positive first exit
at \(y=z/(s|\eta|^2)=\delta_{\rm an}\) satisfies \(\rho_{\rm lo}\le\rho\le\rho_{\rm hi}\).
On this compact interval we first match the two exits, then prepare their
endpoints and the return displacement. The matching equation gives a
rational leading relation with a controlled remainder.
Choose a small fixed reference radius \(v_b>0\), and set
\[
 L_+(\rho)=\log\frac{\rho}{1+\rho},\qquad
 L_-(\rho)=\log\frac{\rho}{1-\rho},\qquad
 \mathfrak c_\sigma(p)=\mathcal A_\sigma(v_b)-L_\sigma(v_b).
\]
On compact positive intervals away from zero, and for
the negative sign away from one,
\begin{equation}
\label{eq:direct-buffered-primary-action}
 \mathcal A_\sigma(\rho)
       =\mathfrak c_\sigma+L_\sigma(\rho)+O_{C^j}(r+s).
\end{equation}
Subtract the integrands on the interval between \(v_b\)
and \(\rho\). Their denominators have fixed positive
margins, and \(R_q(\sigma st)=-1+O_{C^j}(s)\);
integration and differentiation give this estimate.
Equation \eqref{eq:direct-buffered-action-difference} implies
\[
 \Lambda=\exp(\mathfrak c_+-\mathfrak c_-),\qquad
 0<\underline\Lambda_{\rm act}\le\Lambda
     \le\overline\Lambda_{\rm act}<\infty.
\]

The annular height construction extends to any such fixed
compact primary interval with positive drift margin.
For these exits, the increase of action from \(K_0r\)
is at least \(c_1\log(\rho/(K_0r))\).
The equation for the central height therefore makes the expanded
prefix input and its first \(j\) radius derivatives
\(O(s^{-j}\e^{-c_2/s})\), uniformly in small \(r\).
Both signs have thin central prefixes: the positive action
gap can be made larger than their fixed central action
difference by reducing \(r\).
Thus the equation for the central height is
\begin{equation}
\label{eq:direct-buffered-primary-source}
 s\log(h/s^3)=-\mathcal A_\sigma(\rho_\sigma)
                 +sM_\sigma^0(\rho_\sigma),\qquad
 \|M_\sigma^0\|_{C^4}\le C .
\end{equation}
Here \(M_\sigma^0=2\log\rho_\sigma+\log\Psi_\sigma-\log\Upsilon_\sigma\).
Equating the two equations for the central height gives
\[
 L_-(\omega)-L_+(\rho)=\log\Lambda+O_{C^3}(r+s).
\]
For each value of the parameter-dependent multiplier \(\Lambda\),
the leading scalar relation has the solution
\[
 f_\Lambda(\rho)=\frac{\Lambda\rho}{1+(1+\Lambda)\rho}.
\]
For \(\rho_{\rm lo}\le\rho\le\rho_{\rm hi}\), \(f_\Lambda(\rho)\) lies
in a fixed compact subset of \((0,1)\). Comparing the central-height
equations at slightly wider negative exits gives opposite signs;
first-exit order gives the matching negative passage.
The lower bound on \(L_-'\) then gives, by implicit differentiation,
\begin{equation}
\label{eq:direct-buffered-fractional-relation}
 \omega=f_\Lambda(\rho)+O_{C^3}(r+s).
\end{equation}

The rational profile is the inverse, after a sign change and positive
rescaling, of the base-point relation of Huzak--Kristiansen
\cite[(2.13) and (6.8)]{HK2026}.
For this scalar relation, the substitutions
\[
 \rho=-\beta_{\rm HK}x_{\rm out}^{b_{\rm HK}},\qquad
 \omega=\beta_{\rm HK}x_{\rm in}^{b_{\rm HK}},\qquad
 \Lambda=\e^{-K_{\rm HK}},\qquad
 K_{\rm HK}=\frac{\lambda_1^{\rm HK}\pi}
            {\sqrt{-4\lambda_0^{\rm HK}-(\lambda_1^{\rm HK})^2}}
\]
give \(\omega=f_\Lambda(\rho)\); their domain becomes
\(\rho>0\), \(0<\omega<\Lambda/(1+\Lambda)\).
The two-scale \(C^3\) remainder in
\eqref{eq:direct-buffered-fractional-relation} follows by matching the
two central-height equations above.

For the endpoints apply \cref{lem:direct-small-height-continuation}
with \(\ell=\delta=s\), \(x=\rho_\sigma\) and
\(\zeta_0(x)=\delta_{\rm an}x^2\). The coefficient ratios are bounded,
and \(K_\sigma(x)=-s^{-3}q(\sigma sx)=D_\sigma(x)\) has a fixed
positive margin on the compact exit intervals above. The entrance
factor is positive with bounded \(C^3\) norm, including its moving
height. Choose a common \(Z_1=Z_0\) and put
\(b_\sigma^e=b_\sigma\), \(a_\sigma^e=\varpi_{s^2Z_0}A_\sigma\).
Since \(q_2=-s\), the multiplier is
\(\varpi_{s^2Z_0}=1+O(s(1+|\log s|))\). Thus \(a_\sigma^e\) has fixed
positive bounds, and
\begin{equation}
\label{eq:direct-buffered-primary-endpoints}
 I_\sigma(\rho_\sigma)=b_\sigma^e
       -\tfrac12s^2a_\sigma^e\rho_\sigma^2
       +s^2O_{C^3}(s(1+|\log s|)).
\end{equation}

The endpoint preparation now turns the matching relation into an
expansion of the return displacement. Substituting
\eqref{eq:direct-buffered-fractional-relation} and
\eqref{eq:direct-buffered-primary-endpoints} into the exterior map gives
\begin{equation}
\label{eq:direct-buffered-prepared-displacement}
 \frac{2\Delta}{s^2a_+^e}
 =D(p)+\rho^2-\lambda(p)f_\Lambda(\rho)^2
                 +O_{C^3}(r+s(1+|\log s|)),
 \qquad
 \lambda=\Pi_p'(b_-^e)a_-^e/a_+^e .
\end{equation}
Here \(D(p)\) depends only on \(p\), and \(\lambda\) has fixed positive
bounds. For every \(h\in\mathcal H_p\) whose positive first exit lies in
\([\rho_{\rm lo},\rho_{\rm hi}]\), the reference endpoints \(b_\sigma^e\)
differ from the corresponding admissible endpoints by \(O(s^2)\).
Thus, whenever this set of heights is nonempty, the reference endpoints
lie in the enlarged exterior windows. The exterior Taylor remainder adds
\(O_{C^3}(s^2)\) after normalization.
For the principal part \(\widehat\Delta_{\rm m}=D+\rho^2-\lambda f_\Lambda^2\),
\[
 \left(\frac{\widehat\Delta_{\rm m}'}{2\rho}\right)'
 =\frac{3\lambda\Lambda^2(1+\Lambda)}
             {[1+(1+\Lambda)\rho]^4}>c_*>0 .
\]
Apply \cref{lem:direct-derivative-ratio-count} with
\(\phi=\rho^2\), \(\psi=f_\Lambda^2\), and
\(\widehat\Delta=2\Delta/(s^2a_+^e)\). The compact positive radius
interval gives \(\phi'\ge2\rho_{\rm lo}\); the displayed margin
therefore persists under the remainder in
\eqref{eq:direct-buffered-prepared-displacement}. It gives at most two
zeros on this interval and
\[
 \left(\frac{\widehat\Delta_\rho}{2\rho}\right)_\rho>0.
\]
For the crossing direction, use the same endpoint identity as in
\eqref{eq:direct-quadratic-prefix-zero-orientation}, here in the exit
coordinate:
\[
 \frac{\widehat\Delta_\rho}{2\rho}
 =\frac{I_{+,\rho}}{s^2a_+^e\rho}(\e^{\mathcal L}-1),
 \qquad I_{+,\rho}<0.
\]
At a slope zero it gives \(\mathcal L_\rho<0\). Since
\eqref{eq:direct-buffered-primary-source} gives
\(T_\rho=\mathcal A_+'-s(M_+^0)'>0\), every such zero has
\(\mathcal L_T<0\).

\subsubsection{Large positive exits}

We now treat the central heights in \(\mathcal H_p\) whose positive
half-orbit is thin at \(\rho_i\), with the prefix
\eqref{eq:direct-annulus-prefix}, and whose adaptive exit, at
\(z=\delta_{\rm an}\varrho^2(s+\varrho)\) in the notation of
\cref{lem:positive-branch-adaptive-exit}, has radius
\(\varrho_e=s\rho_e\) with \(\rho_e\ge1\). The adaptive exit is the
unique upward crossing of that level, so \(\rho_e\ge1\) means
\(z(s;h)\le2\delta_{\rm an}s^3\). By \cref{lem:central-height-order},
this condition and thinness at \(\rho_i\) each define an interval;
their intersection is the interval considered here. For each such height,
the negative annular thin exit at \(z/(s|\eta|^2)=\delta_{\rm an}\) remains
in a fixed compact interval strictly below its far root, and for large
\(\rho_e\) the return contracts.

At \(\rho_i\), \eqref{eq:direct-annulus-prefix} reads
\[
 K_0^2\Psi_+(1/K_0,K_0r,y_i)=\xi_+\Upsilon_+(\xi_+),\qquad
 y_i=\frac{z(\varrho_i;h)}{K_0^2Q},\qquad \varrho_i=K_0sr .
\]
Inverting this height coordinate gives
\[
 V_{\rm an}(\xi)=\frac{K_0^2}{\xi}
   \Psi_+^{-1}\bigl(1/K_0,K_0r;\xi\Upsilon_+(\xi)/K_0^2\bigr),
 \qquad z(\varrho_i;h)/Q=\xi_+V_{\rm an}(\xi_+).
\]
Thinness keeps the inverse in its fixed domain. Its derivative at zero
is one, so \(V_{\rm an}(0)=1\), and the annular coordinate bounds
supply the unit, reciprocal and inverse bounds required by
\eqref{eq:positive-adaptive-prefix}. Thus
\cref{lem:positive-branch-adaptive-exit} applies with
\(\mathcal V=V_{\rm an}\) and \(\mathcal A_i=\mathcal A_{+,i}\).

For \(1\le\rho\le\varrho_*/s\), where \(\varrho_*\) is a small
fixed unscaled radius, subtraction of the positive
action integrands gives the value estimate
\[
 \mathcal A_+(\rho)=\mathfrak c_++L_+(\rho)
            +O\{r+s(1+\log(1+\rho))\}.
\]
Indeed \(D_+(t)\asymp t^2(1+t)\) and
\(R_q(st)+1=O(s+st)\); the difference is bounded by
\(Cr/t^2+Cs/(1+t)\) for \(t\ge v_b\).
The action increment from \(K_0r\) to \(1\) gives the gap in part
(iv) of the lemma, with \(\rho_{\rm g}=1\): since \(D_+>0\), the annular
bound \(\mathcal J_+(\rho)\ge c\log(\rho/(K_0r))\) at \(\rho=v_1\) gives
\(\mathcal J_+(1)\ge \mathcal J_+(v_1)\ge\frac c2\log(1/r)\) for small \(r\).
Taking logarithms in \eqref{eq:direct-no-root-adaptive-input}, with
\(\mathcal A_{+,i}+s\mathcal J_a(s\rho_e)=\mathcal A_+(\rho_e)\) and
\(\varrho_e^2(s+\varrho_e)/Q=\rho_e^2(1+\rho_e)/r^2\), gives the equation
for the central height at the adaptive exit:
\[
 s\log(h/s^3)=-\mathcal A_+(\rho_e)
   +s\log[\rho_e^2(1+\rho_e)\Psi_a(u_e,v_e,\varrho_e,\delta_{\rm an})]
   -s\log\Upsilon_a(\xi_+).
\]
By \eqref{eq:positive-adaptive-input-weight}, \(\xi_+\) and
\(\log\Upsilon_a(\xi_+)\) are at most \(Cs^{-3}r^{c/(2s)-2}\). After
reducing \(r\), the negative input
\(\xi_-=\e^{(\mathcal A_{-,i}-\mathcal A_{+,i})/s}\xi_+\) is exponentially
small as well, since
\eqref{eq:direct-buffered-action-difference} bounds the difference of the
central actions. The remaining height prefactor
\(\rho_e^2(1+\rho_e)\Psi_a\) is at least \(c\delta_{\rm an}\) and, since the
unscaled exit radius is bounded, at most a fixed power of \(s^{-1}\).
Thus
\[
 s\log(h/s^3)=-\mathfrak c_+-L_+(\rho_e)
                    +O(r+s(1+|\log s|)).
\]
Since \(f_\Lambda(\rho_e)\), for \(\rho_e\ge1\), lies
in a fixed compact subset of \((0,1)\), comparison with
\eqref{eq:direct-buffered-primary-source} at two wider
negative endpoints proves existence of the negative exit.
Keep these negative comparison endpoints fixed when choosing the
large-exit threshold below.

Differentiation of the equation for the central height on the
negative branch gives
\[
 \frac{d\omega}{d\log h}
   =-\frac{s}{\mathcal A_-'(\omega)
                         -s(M_-^0)'(\omega)},\qquad
 0<hI_{-,h}\le Cs^3
\]
by \eqref{eq:direct-buffered-primary-endpoints}.
On the positive branch, parts (ii) and (iv) of
\cref{lem:positive-branch-adaptive-exit} give
\(d\lambda_{\rm in}/d\log h=1+O(\xi_+)=1+O(s^N)\). By part (iii), the
landing satisfies \eqref{eq:direct-tail-smallness}. These are the
hypotheses of \eqref{eq:direct-adaptive-source-slope}. Since \(q_2=-s\),
the condition \(q_2\le C_2s\) of \cref{cor:large-exit-contraction} holds
with \(C_2=0\); take \(C_1\) the constant of the negative bound above, and
fix \(R_0\ge1\) satisfying the condition on \(M_1\) there. The corollary
gives \(\mathcal L<0\), and at most one zero of \(\Delta_p\), for
\(\rho_e\ge R_0\).
Writing \(\rho\) for the radius of the annular thin exit and \(\rho_e\)
for that of the adaptive exit, we have
\(\rho\le\rho_e\) by \eqref{eq:positive-annular-adaptive-heights}.
On the interval \(\rho\ge\rho_{\rm lo}>0\),
any zero of \(\mathcal L\) must have \(\rho_e<R_0\), because
the large-exit estimate gives \(\mathcal L<0\) otherwise.
Such a central height has \(\rho_{\rm lo}\le\rho<R_0\), so the compact-primary
comparison proves \(\mathcal L_T<0\) at that zero. Thus every
zero of \(\mathcal L\) on this entire later interval of central
heights is a strict downward crossing.

\subsubsection{Covering all central heights}

We now join the comparisons in their central-height order. On the
overlaps they concern the same slope ratio, so their crossing
directions give one zero count on the whole interval.
Use the same sufficiently small thin threshold \(\delta_{\rm an}\)
at \(z_e=s^3r^2\delta_{\rm an}\). The quadratic-prefix lemma,
with \(\delta\) there replaced by \(s\), \(\ell=sr\),
and comparison drifts \(a-\sigma bx+x^2\), gives at most two
zeros on
\[
 h\ge z_e,\qquad\hbox{or}\qquad
 h<z_e,\quad |\eta_+(z_e;h)|/(sr)\le M_0 .
\]
The compact negative inverse follows from the first assertion of
\cref{lem:prefix-root-cover},
since \(a\ge\nu/4\) and the positive margin holds up to \(M_0\).
Every other central height has its positive half-orbit below
height \(z_e\) at all radii up to \(M_0sr\). For \(K_0r\le\rho\le M_0r\),
it therefore satisfies
\[
 y(\rho)<\delta_{\rm an}(r/\rho)^2<\delta_{\rm an}.
\]
Its positive thin exit occurs beyond \(M_0r\).
Such an exit exists before the admissible endpoint: the
boundary coefficient satisfies \(q(\eta)<0\) on the positive branch throughout the
small box in the coordinates \((\eta,z)\), whereas the endpoint has
\(y\asymp s^{-1}>\delta_{\rm an}\). At a crossing of \(y=\delta_{\rm an}\),
the right side of \eqref{eq:direct-annulus-level-crossing} is positive,
because \(sK_+=O(s+|\eta|)\) past the inner cut.
Thus every crossing is upward, and the first exit is unique.

Take \(\rho_{\rm lo}=v_0/2\) in the compact-primary comparison and reduce \(r\)
so that \(M_0r<\rho_{\rm lo}\). In increasing
\(T\), the early central heights, the annular interval
\(M_0r\le\rho\le v_0\), and the later interval \(\rho\ge\rho_{\rm lo}\)
can be cut within their overlaps to give an ordered cover of
\(\mathcal H_p\). Their order
follows from the decrease of the central height with the exit radius.
The first join has an overlap: an annular exit at \(M_0r\)
has height \(s^3r^2M_0^2\delta_{\rm an}>z_e\), whereas at the height
threshold of the early interval the positive half-orbit has height
\(z_e\) at that radius. The second overlap is
\(v_0/2\le\rho\le v_0\). These assertions persist after
intersection with \(\mathcal H_p\), including empty end pieces.

On the early interval, \cref{lem:direct-quadratic-prefix} gives
\(\mathcal L_T>0\) at every zero of \(\mathcal L\).
On the middle interval, either \(\mathcal L''<0\), for
\(b\ge b_{*,-}\), or \(\mathcal L<0\), for \(b\le b_{*,-}\).
In the latter case merge the nonvanishing middle interval with the
initial interval; all zeros on their union cross upwards.
On the later interval every zero has \(\mathcal L_T<0\), by the
compact-primary and large-exit comparisons above.
\Cref{lem:direct-oriented-zero-count} therefore gives at most two
zeros of \(\mathcal L\) on the whole interval \(\mathcal H_p\), and
\cref{cor:ordered-cover-zero-count} gives at most three zeros of
\(\Delta_p\). This proves \cref{prop:direct-buffered-negative}.
